\documentclass[11pt]{article}\usepackage{amsmath,amssymb,amsthm,geometry}\geometry{margin=1in}
\newtheorem{lemma}{Lemma}
\title{The full-grade target for iterated continuation repair}\date{16 November 2026}\begin{document}\maketitle
\paragraph{Scope.} Exact grade lemma and next-runtime obligation, not an iterated stochastic assembly theorem. Work in the ordered polynomial collection doctrine with tagged coproduct, full all-small shape families and actual tensor coefficients. No unused shapes are removed. $I$ is the unary tensor-unit collection; $0$ is the empty collection. The recursive color set $U$ is fixed.

For continuation grades $D,E$, there is a canonical full collection arrow
\[
 j_{D,E}:(D\circ E)+I\longrightarrow(D+I)\circ(E+I).       \tag{1}
\]
On a shape in $D\circ E$, choose the $D$ branch at its root and the $E$ branch at every child position. On the separate $I$ shape, choose the $I$ branch at both layers. The shape map is injective because branch tags distinguish these cases. Ordered external positions are unchanged. The actual coefficient arrow is the labeled tensor associator/unitor, an isomorphism at each source shape. This defines (1) on every source shape, including unused shapes and all coefficients.

\begin{lemma}
For arbitrary $C,D,E$, the following full terminal collections map into one common target:
\[
\begin{split}
 &(R_U(C\circ(D+I)))\circ E
       \longrightarrow R_U(C\circ(D+I)\circ(E+I)),\\
 &R_U(C\circ((D\circ E)+I))
       \longrightarrow R_U(C\circ(D+I)\circ(E+I)).       \tag{2}
\end{split}
\]
Both arrows are injective on full shapes, preserve ordered external positions and are isomorphisms on each source-shape coefficient. Neither is claimed surjective.
\end{lemma}
\begin{proof}
The first is the established whole-case post-tightening grade map, applied to $C\circ(D+I)$ and $E$. Each existing internal feedback edge takes the new $I$ decoration and each external decoration takes its supplied $E$ shape. The second is $R_U(C\circ j_{D,E})$, using the whole finite-tree functor. At every node use the tagged choice in (1), retaining each loop label and every existing edge. Finite-tree recursion maps all shapes and actual tensor coefficients, with no reference to realizing states. Injection and arity preservation follow by reconstructing the original tags/edges. The coefficient maps are composites of the specified tensor isomorphisms.
\end{proof}

A reverse whole-grade map cannot be presumed. Take $C=I$, $D=0$ and $E$ with one binary unit-coefficient shape. Then the smaller root grade $C\circ((D\circ E)+I)$ is unary $I$, whereas the target root grade is $E+I$ and has a binary shape. With a singleton $U$, every finite shape of $R_UI$ has arity one, whereas $R_U(E+I)$ has a binary leaf shape. No arity-preserving map from the latter to the former exists. This is an unused-grade obstruction even if no completed $Q$ state realizes a $D$ body.

\paragraph{Associativity of the specified embeddings.}
For three continuation layers $D,E,F$, both ways of embedding the genuine composite branch $((D\circ E)\circ F)+I\cong(D\circ(E\circ F))+I$ into
\[
                  (D+I)\circ(E+I)\circ(F+I)
\]
make every genuine layer active and every identity bypass take its $I$ branch at every layer. Thus they give the same tagged shape and ordered position map. The coefficient maps agree by the labeled tensor associativity/unit coherence, not by faithfulness of a scalar body. The same equation holds after substitution into $C$ and finite-tree application of $R_U$. This is a precise grade equation; it is not a proof that stochastic runtime comparisons have the same domains or preserve their law families.

\paragraph{Independent runtime still required.}
A three-tier target should carry primitive $P$ invocation trees, $Q$ copies at their $B$ ports and $S$ copies at the resulting $Z$ ports. An $S$ request may depend on its producing $Q$ initializer witness, in addition to the producing $P$ ancestry. It must not condition on unrelated $Q$ siblings or remote paid atoms. All silent descendants retain the paid owner of their first request; old paid $Q$ actors can independently own new $S$ requests. Initial coefficients remain pure tensors and paid coefficients retain whole owner blocks. The natural maximal full grade is the common target in (2), with actual bodies selecting only the typed cases. Whether all generated stochastic routes map there and commute requires initial/paid law coarsening, actual coefficient/body comparisons and full unused writer maps. None is inferred solely from (1)--(2). A smaller colored runtime may have a grade isomorphism, but is a different doctrine.
\end{document}
